%% =====================================================================
\chapter{Backend Foundation}\label{ch:foundation}

Before any feature, every request passes through the same foundation: the
configuration, the Express application with its security middleware, input
validation and the central error handler. This chapter explains each piece.

\section{Life of a request}

Take \code{GET \codesl{}api\codesl{}records?patientId=...} sent by a doctor. In order:

\begin{enumerate}
  \item \textbf{helmet} adds security headers (CSP, \code{nosniff}, \ldots).
  \item \textbf{CORS}, \textbf{cookie-parser} and the \textbf{JSON parser}
    (100\,kB limit) run.
  \item The \textbf{rate limiter} counts the request against the caller's IP.
  \item The router for \code{\codesl{}api\codesl{}records} runs \textbf{authenticate}: verify
    the JWT, check it is not revoked, load the account state, compute the
    trust score (Chapter~\ref{ch:zerotrust}).
  \item \textbf{requireRole}\code{('patient','doctor', \{minTrustScore:70, allowBreakGlass:true\})}
    checks the role, approval status and trust score.
  \item \textbf{validate} checks the query string against a Joi schema.
  \item The handler calls \textbf{consentService} (does this doctor have
    consent?) and then MongoDB.
  \item If anything throws, the \textbf{error handler} turns it into JSON.
  \item Allow and deny decisions are written to the \textbf{access chain} of
    the ledger.
\end{enumerate}

\section{Configuration: \file{src\codesl{}config\codesl{}config\codedot{}js}}

\textbf{Purpose.} Turn environment variables into one typed object, and refuse
to start if the configuration is missing or unsafe.

\textbf{Why.} The original service started with missing secrets and failed
later on the first login, with an obscure error. Failing at start-up is cheaper
to debug and can never half-work in production.

\begin{lstlisting}[language=JavaScript,caption={Start-up validation (\file{src/config/config.js})}]
function validateConfig(cfg = config, rawEnv = env) {
  const required = ['MONGODB_URI', 'JWT_SECRET', 'JWT_REFRESH_SECRET'];
  if (cfg.email.transport === 'smtp') {
    required.push('EMAIL_USER', 'EMAIL_PASSWORD');
  }

  const missing = required.filter((key) => !rawEnv[key] || String(rawEnv[key]).trim() === '');
  if (missing.length > 0) {
    throw new Error(`Missing required environment variables: ${missing.join(', ')}`);
  }

  if (!['smtp', 'console', 'json'].includes(cfg.email.transport)) {
    throw new Error(`EMAIL_TRANSPORT must be smtp, console or json (got "${cfg.email.transport}")`);
  }

  // Printing OTPs to a log is only acceptable on a developer machine.
  if (cfg.isProduction && cfg.email.transport !== 'smtp') {
    throw new Error('EMAIL_TRANSPORT must be "smtp" in production');
  }

  if (rawEnv.JWT_SECRET === rawEnv.JWT_REFRESH_SECRET) {
    throw new Error('JWT_SECRET and JWT_REFRESH_SECRET must be different');
  }

  return true;
}
\end{lstlisting}

\begin{howto}
\begin{itemize}
  \item Email credentials are required only when mail is really sent
    (\code{smtp}).
  \item The \code{console} transport prints OTP codes to the log, which is
    convenient on a laptop and a security hole on a server, so production
    rejects it.
  \item If both JWT secrets were equal, a refresh token could be presented as an
    access token. Requiring different secrets closes that.
  \item The function takes \code{cfg} and \code{rawEnv} as parameters so unit
    tests can call it with fake environments.
\end{itemize}
\end{howto}

\section{Process start-up: \file{src\codesl{}server\codedot{}js}}

\begin{lstlisting}[language=JavaScript,caption={Start-up order (\file{src/server.js})}]
async function start() {
  try {
    config.validateConfig();
    await connectMongo();
    await connectRedis();
    await HospitalSettings.getSingleton();
    await emailService.verifyConnection();

    const app = require('./app');
    const server = app.listen(config.port, () => {
      logger.info(`MediChain API listening on http://localhost:${config.port} (${config.nodeEnv})`);
    });

    const shutdown = (signal) => {
      logger.info(`${signal} received, shutting down`);
      server.close(async () => {
        await Promise.allSettled([disconnectMongo(), disconnectRedis()]);
        process.exit(0);
      });
      setTimeout(() => process.exit(1), 10000).unref();
    };
    process.on('SIGTERM', () => shutdown('SIGTERM'));
    process.on('SIGINT', () => shutdown('SIGINT'));
  } catch (error) {
    logger.error(`Failed to start: ${error.message}`);
    process.exit(1);
  }
}
\end{lstlisting}

\begin{howto}
The server listens only after every dependency answered. \code{connectRedis}
sends a \code{PING} and throws if Redis is not there (the old service swallowed
this error and then failed every login). \code{HospitalSettings.getSingleton()}
creates the settings document on first run. \code{verifyConnection} logs into
SMTP so a wrong Gmail app password shows up in the log at start-up, not at the
first login. On \code{SIGTERM} (Docker stop) the server stops accepting
connections, finishes in-flight requests and closes the databases; if that takes
longer than 10\,s it exits anyway.
\end{howto}

\section{The Express application: \file{src\codesl{}app\codedot{}js}}

\subsection{Security headers and the Content-Security-Policy}

\begin{lstlisting}[language=JavaScript,caption={helmet and CSP (\file{src/app.js})}]
// The frontend is served by this same server, so a strict Content-Security-Policy
// can allow scripts only from this origin plus two pinned CDNs. Inline <script>
// and onclick="" handlers are blocked, which shuts the main door to XSS.
app.use(helmet({
  contentSecurityPolicy: {
    useDefaults: true,
    directives: {
      'script-src': ["'self'", 'https://cdn.jsdelivr.net', 'https://unpkg.com'],
      'style-src': ["'self'", "'unsafe-inline'", 'https://fonts.googleapis.com', 'https://cdn.jsdelivr.net'],
      'font-src': ["'self'", 'https://fonts.gstatic.com', 'data:'],
      'img-src': ["'self'", 'data:', 'blob:', 'https://images.pexels.com'],
      'connect-src': ["'self'"],
      'frame-src': ["'self'", 'blob:'],
      'object-src': ["'self'", 'blob:'],
      'upgrade-insecure-requests': config.isProduction ? [] : null,
    },
  },
  crossOriginEmbedderPolicy: false,
}));
\end{lstlisting}

\begin{howto}
A Content-Security-Policy tells the browser which sources it may load code from.
With \code{script-src 'self'} and no \code{'unsafe-inline'}, even if an attacker
manages to get \code{<script>alert(1)<\codesl{}script>} into a patient's name, the browser
refuses to run it. This is why the frontend has no inline scripts and no
\code{onclick=} attributes: all handlers are attached from \code{.js} files.
\code{connect-src 'self'} means page scripts can only call this API. \code{blob:}
in \code{frame-src}/\code{object-src} lets the dashboards preview downloaded PDFs.
\end{howto}

\subsection{Rate limits}

\begin{lstlisting}[language=JavaScript,caption={Two rate limiters (\file{src/app.js})}]
app.use('/api/', rateLimit({
  windowMs, max: config.rateLimit.max, standardHeaders: true, legacyHeaders: false,
  message: limiterResponse('Too many requests from this IP, please try again later.'),
  skip: () => config.isTest,
}));

// Stricter budget for the endpoints an attacker would hammer to guess passwords or codes.
const authLimiter = rateLimit({
  windowMs, max: config.rateLimit.authMax, standardHeaders: true, legacyHeaders: false,
  message: limiterResponse('Too many login attempts from this IP. Try again in 15 minutes.'),
  skip: () => config.isTest,
});
app.use(['/api/auth/login', '/api/auth/register'], authLimiter);
\end{lstlisting}

The general limit protects the whole API; the second, much smaller budget
(30 per 15 minutes by default) applies to login and registration, the endpoints
used for password and OTP guessing. Per-account limits (lockout, OTP attempts)
are added on top in Chapter~\ref{ch:auth}, so an attacker spreading requests
over many IPs still hits a wall.

\subsection{Routes and the bundled frontend}

\begin{lstlisting}[language=JavaScript,caption={Module mounting and static frontend (\file{src/app.js})}]
app.use('/api/auth', require('./modules/auth/auth.routes'));
app.use('/api', require('./modules/users/users.routes'));
app.use('/api/appointments', require('./modules/appointments/appointments.routes'));
app.use('/api/consents', require('./modules/consents/consents.routes'));
app.use('/api/records', require('./modules/records/records.routes'));
app.use('/api/files', require('./modules/records/files.routes'));
app.use('/api/prescriptions', require('./modules/prescriptions/prescriptions.routes'));
app.use('/api', require('./modules/vitals/vitals.routes'));
app.use('/api/bills', require('./modules/billing/billing.routes'));
app.use('/api/claims', require('./modules/claims/claims.routes'));
app.use('/api/notifications', require('./modules/notifications/notifications.routes'));
app.use('/api/admin', require('./modules/admin/admin.routes'));
app.use('/api/ledger', require('./modules/ledger/ledger.routes'));
app.use('/api', notFoundHandler);

// Same origin as the API: the SameSite=Strict refresh cookie just works and no
// CORS is needed in the browser. Development files are never served.
app.use((req, res, next) => (/^\/(node_modules|tests|package(-lock)?\.json|jest\.config)/.test(req.path) ? notFoundHandler(req, res) : next()));
app.use(express.static(config.frontendDir, { index: 'homepage.html', extensions: ['html'] }));
\end{lstlisting}

The \code{\codesl{}api} 404 handler is mounted \emph{before} the static files, so a
mistyped API path returns a JSON 404 instead of an HTML page. The regular
expression blocks the frontend's development files (its \code{node\_modules},
tests and \code{package.json}) from being served.

\section{Errors: \file{src\codesl{}utils\codesl{}errors\codedot{}js} and the error handler}

\begin{lstlisting}[language=JavaScript,caption={AppError and asyncHandler (\file{src/utils/errors.js})}]
class AppError extends Error {
  constructor(status, message, code, details) {
    super(message);
    this.name = 'AppError';
    this.status = status;
    this.code = code || undefined;
    this.details = details || undefined;
  }
}

const badRequest = (message, details) => new AppError(400, message, 'BAD_REQUEST', details);
const unauthorized = (message, code = 'UNAUTHORIZED') => new AppError(401, message, code);
const forbidden = (message, code = 'FORBIDDEN') => new AppError(403, message, code);
const notFound = (message = 'Resource not found') => new AppError(404, message, 'NOT_FOUND');
const conflict = (message, code = 'CONFLICT') => new AppError(409, message, code);
// ...
const asyncHandler = (fn) => (req, res, next) => Promise.resolve(fn(req, res, next)).catch(next);
\end{lstlisting}

\begin{howto}
Handlers simply \code{throw conflict('This slot has just been booked...', 'SLOT\_TAKEN')}.
Express 4 does not catch errors thrown from \code{async} functions, so every
handler is wrapped in \code{asyncHandler}, which forwards a rejected promise to
\code{next(error)}. The central handler then produces the response. The
machine-readable \code{code} lets the frontend react to specific errors (for
example \code{SESSION\_ENDED} sends the user back to the login page).
\end{howto}

\begin{lstlisting}[language=JavaScript,caption={Central error handler, mapping section (\file{src/middleware/errorHandler.js})}]
  if (err.message === 'REDIS_UNAVAILABLE') {
    status = 503;
    message = 'Security service temporarily unavailable. Please try again.';
    code = 'SERVICE_UNAVAILABLE';
  } else if (err instanceof multer.MulterError) {
    status = err.code === 'LIMIT_FILE_SIZE' ? 413 : 400;
    message = err.code === 'LIMIT_FILE_SIZE' ? 'File is too large' : err.message;
    code = err.code;
  } else if (err.name === 'ValidationError') {
    status = 400;
    message = Object.values(err.errors).map((e) => e.message).join(', ');
  } else if (err.name === 'CastError') {
    status = 400;
    message = `Invalid ${err.path}`;
  } else if (err.code === 11000) {
    status = 409;
    message = 'A record with these details already exists';
    code = 'DUPLICATE';
  } else if (err.type === 'entity.too.large') {
    status = 413;
    message = 'Request body is too large';
  } else if (err.type === 'entity.parse.failed') {
    status = 400;
    message = 'Malformed JSON body';
  }

  if (status >= 500) {
    logger.error(`${req.method} ${req.originalUrl} -> ${err.stack || err.message}`);
    if (config.isProduction) message = 'Internal server error';
  }
\end{lstlisting}

Each branch converts a low-level failure into a meaningful HTTP status. The most
important one is the first: when Redis is down, the security layer cannot check
OTPs, refresh tokens or revocations, so the request fails with \textbf{503}
(``fail closed'') rather than being let through. In production, details of
unexpected (500) errors are logged but not sent to the client.

\section{Input validation: \file{src\codesl{}middleware\codesl{}validate\codedot{}js}}

\begin{lstlisting}[language=JavaScript,caption={The validate middleware (\file{src/middleware/validate.js})}]
const validate = (schemas) => (req, res, next) => {
  const errors = [];
  for (const part of ['params', 'query', 'body']) {
    if (!schemas[part]) continue;
    const { error, value } = schemas[part].validate(req[part] || {}, {
      abortEarly: false, stripUnknown: false, convert: true,
    });
    if (error) {
      errors.push(...error.details.map((d) => d.message.replace(/"/g, '')));
    } else if (part === 'query') {
      Object.defineProperty(req, 'query', { value, writable: true, configurable: true });
    } else {
      req[part] = value;
    }
  }
  if (errors.length) {
    return res.status(400).json({ success: false, message: 'Validation failed', errors });
  }
  return next();
};

const objectId = () => Joi.string().pattern(/^[a-f\d]{24}$/i).message('must be a valid id');

const password = () => Joi.string()
  .min(8)
  .max(128)
  .pattern(/^(?=.*[a-z])(?=.*[A-Z])(?=.*\d)(?=.*[@$!%*?&#^()\-_=+]).+$/)
\end{lstlisting}

\begin{howto}
Each route declares what it accepts, for example
\code{validate(\{ body: Joi.object(\{ doctorId: objectId().required(), ... \}) \})}.
\begin{itemize}
  \item \code{abortEarly:false} reports \emph{every} problem at once, so a form
    can show all errors together.
  \item \code{stripUnknown:false} means unknown fields are an error, not silently
    dropped: a client cannot sneak in \code{role:'admin'}.
  \item \code{convert:true} turns \code{"7"} into \code{7} and ISO strings into
    dates; the cleaned value replaces the original, so handlers work with
    trusted, typed data. (\code{req.query} is a getter in newer Express, hence
    \code{defineProperty}.)
\end{itemize}
\end{howto}

\section{Databases: \file{src\codesl{}db\codesl{}}}

\begin{lstlisting}[language=JavaScript,caption={GridFS bucket (\file{src/db/mongo.js})}]
/**
 * GridFS stores files larger than a normal document by splitting them into
 * 255 kB chunks (collections `uploads.files` and `uploads.chunks`). It ships
 * with the MongoDB driver, so no extra infrastructure is needed (decision D3).
 */
function getBucket() {
  if (!bucket) {
    bucket = new mongoose.mongo.GridFSBucket(mongoose.connection.db, { bucketName: 'uploads' });
  }
  return bucket;
}
\end{lstlisting}

\begin{lstlisting}[language=JavaScript,caption={Redis connection that fails loudly (\file{src/db/redis.js})}]
async function connectRedis() {
  if (!redisClient.isOpen) {
    await redisClient.connect();
  }
  await redisClient.ping();
  return redisClient;
}
\end{lstlisting}

\section{Hashing helpers: \file{src\codesl{}utils\codesl{}hash\codedot{}js}}

These three functions underpin the ledger (Chapter~\ref{ch:ledger}).

\begin{lstlisting}[language=JavaScript,caption={Canonical JSON and hashing (\file{src/utils/hash.js})}]
const sha256 = (value) => crypto.createHash('sha256').update(typeof value === 'string' || Buffer.isBuffer(value) ? value : String(value ?? '')).digest('hex');

/**
 * Deterministic JSON: object keys are sorted at every level, so the same data
 * always produces the same string (and therefore the same hash) no matter in
 * which order the fields were assigned.
 */
function canonicalJSON(value) {
  if (value === null || value === undefined) return 'null';
  if (value instanceof Date) return JSON.stringify(value.toISOString());
  if (Array.isArray(value)) return `[${value.map(canonicalJSON).join(',')}]`;
  if (typeof value === 'object') {
    if (typeof value.toHexString === 'function') return JSON.stringify(value.toHexString()); // ObjectId
    const keys = Object.keys(value).filter((k) => value[k] !== undefined).sort();
    return `{${keys.map((k) => `${JSON.stringify(k)}:${canonicalJSON(value[k])}`).join(',')}}`;
  }
  return JSON.stringify(value);
}

const hashObject = (value) => sha256(canonicalJSON(value));
\end{lstlisting}

\begin{howto}[Why canonical JSON matters]
\code{JSON.stringify(\{a:1,b:2\})} and \code{JSON.stringify(\{b:2,a:1\})} produce
different strings, and therefore different hashes, for the same data. MongoDB
does not promise to return fields in the order they were written. Sorting the
keys at every level, and writing dates and ObjectIds in one fixed format, makes
the hash a function of the \emph{data} only. Without this, verification would
report false tampering.
\end{howto}
